\documentclass[11pt]{article}
% Shared preamble for the hanzo-kernel Paper 07 series.
% One style, one vocabulary, inputted by every paper so the series reads as one voice.
\usepackage[margin=1in]{geometry}
\usepackage{booktabs}
\usepackage{amsmath}
\usepackage{amssymb}
\usepackage{graphicx}
\usepackage{xcolor}
\usepackage{hyperref}
\hypersetup{colorlinks=true, linkcolor=blue, urlcolor=blue, citecolor=blue}
\usepackage{enumitem}
\usepackage{listings}
\usepackage{array}

\definecolor{kw}{rgb}{0.13,0.29,0.53}
\definecolor{cmt}{rgb}{0.40,0.40,0.40}
\definecolor{str}{rgb}{0.16,0.50,0.16}
\lstset{
  basicstyle=\ttfamily\footnotesize,
  keywordstyle=\color{kw}\bfseries,
  commentstyle=\color{cmt}\itshape,
  stringstyle=\color{str},
  showstringspaces=false,
  breaklines=true,
  columns=fullflexible,
  frame=single,
  framesep=4pt,
}
\lstdefinelanguage{Rust}{
  morekeywords={fn,let,mut,pub,struct,impl,trait,enum,match,if,else,for,while,loop,return,use,mod,crate,self,super,async,await,where,type,const,static,unsafe,ref,move,dyn,as,true,false},
  morecomment=[l]{//}, morecomment=[s]{/*}{*/}, morestring=[b]", sensitive=true,
}

\newcommand{\x}{\ensuremath{\times}}
\newcommand{\dpa}{\texttt{dp4a}}
\newcommand{\hk}{\texttt{hanzo-kernel}}
\newcommand{\code}[1]{\texttt{#1}}

% Absorb minor prose overfulls without rivers; keep line-breaking sane.
\tolerance=800
\emergencystretch=3em


\title{\textbf{The Schedule Is a Value:\\ Comptime-Specialized Autotuning Without Forks}}
\author{Hanzo AI --- ML Systems}
\date{hanzo-kernel Paper 07.2 --- part of the \emph{One Source, Every Backend} series}

\begin{document}
\maketitle

\begin{abstract}
A single-source kernel DSL promises to kill the ``one operation, $N$ backends, $N$ numeric behaviors''
bug class by writing each op once. But there is a subtler place the same bug hides: the \emph{schedule}
--- the tile size, vector width, unroll factor, and thread-block dimension a kernel is tuned with, which
differs per device. The usual fix is to fork a hand-tuned kernel per device, and a fork is exactly what
reintroduces the bug: two schedules, drifting independently, one of them one day reordering an
accumulation and changing the numbers. This paper shows how to autotune without forking. The schedule is
exposed as a set of compile-time (\code{\#[comptime]}) knobs on the \emph{one} source; each knob value is
a distinct compiled kernel the lowering engine monomorphizes; a small set of these variants is timed once
per \code{(device, op, shape)} and the winner cached to disk. Because the knobs never touch the
arithmetic, \emph{every variant is byte-identical} --- the tuner buys speed and can never spend
correctness. We give the three-value algebra (\code{Variant}, \code{Tuned}, \code{Tuner}), the
per-device on-disk cache, the two shipped variant tables, and the offline tests that prove
select/cache/reload works without a GPU.
\end{abstract}

\section{The pain: a fork that drifts}
\label{sec:pain}

Suppose you have written a matvec kernel once, in a portable DSL, and it lowers correctly to every
backend. You are not done, because ``correct everywhere'' is not ``fast everywhere.'' The fastest
schedule for an AMD iGPU is not the fastest for an NVIDIA datacenter part: they differ in warp/subgroup
width, in how much shared memory pays off, in the ideal number of output elements a thread should carry
(the unroll factor), in the block size that fills the machine. The reflex is to fork --- keep a
\code{matvec\_amd} and a \code{matvec\_nvidia}, each hand-tuned.

That reflex is the original sin the DSL exists to prevent, wearing a schedule's clothes. Two forks are
two kernels, maintained by two hands, drifting apart. The day someone reorders a loop in one fork ``to
help the unroll,'' its accumulation order changes, its rounding changes, and now the same operation
produces different bits on two backends --- the precise failure a single-source DSL was supposed to make
impossible (\cite{oracle}). Worse, the drift is invisible until a model degenerates on one device and not
another.

We want the performance a per-device schedule buys \emph{without} the second kernel that buys the bug.

\section{The idea: a schedule is a compile-time value}
\label{sec:idea}

The DSL already gives one source lowered to every backend. The remaining fork --- the schedule --- can be
folded into that one source if we treat the schedule not as \emph{code} but as a \emph{value}. In \hk{},
a kernel parameter marked \code{\#[comptime]} is resolved at compile time, and the lowering engine
\emph{monomorphizes} the kernel on it: each distinct comptime value produces a distinct compiled kernel,
specialized as if you had hand-written it. So we lift the schedule knobs to \code{\#[comptime]}
parameters of the one kernel (Listing~\ref{lst:tunedkernel}). A ``variant'' is then not a forked file but
a \emph{choice of comptime tuple}. There is one source; only the chosen tuple differs per device.

\begin{lstlisting}[language=Rust,caption={One autotuned source. The unroll factor \code{rpt}
(rows-per-thread) is a comptime knob; the block size is the host's knob. Every \code{(rpt, block)} is a
distinct compiled kernel yet bit-identical to the reference: the per-row fold is the same order, so only
the thread$\to$row mapping and occupancy differ.},label={lst:tunedkernel}]
#[kernel(targets(cuda, metal, vulkan, webgpu, cpu), unchecked)]
pub fn rms_norm_tuned<F: Float>(x: &Array<F>, w: &Array<F>, out: &mut Array<F>,
                                eps: &Array<F>, #[comptime] n: usize, #[comptime] rpt: usize) {
    let nrows = out.len() / n;
    let base = ABSOLUTE_POS * rpt;
    #[unroll]
    for r in 0..rpt {
        let row = base + r;
        if row < nrows {
            let b = row * n;
            let mut ss = F::new(0.0);
            for i in 0..n { let v = x[b + i]; ss += v * v; }   // same fold, same order, every variant
            let denom = (ss / F::cast_from(n as u32) + eps[0]).sqrt();
            for i in 0..n { out[b + i] = x[b + i] / denom * w[i]; }
        }
    }
}
\end{lstlisting}

The critical invariant is stated in the comment and enforced by test: the knobs change the
thread-to-work mapping and the occupancy, \emph{not} the arithmetic. The per-row sum-of-squares is the
same left fold in the same order for every \code{(rpt, block)}. Therefore all variants are
\textbf{byte-identical} to each other and to the reference. The autotuner is choosing among numerically
indistinguishable kernels; it cannot pick a ``faster but slightly wrong'' one, because there is no wrong
one to pick.

\section{The mechanism: three orthogonal values}
\label{sec:mech}

Autotuning decomposes into three plain values, each testable in isolation (Listing~\ref{lst:tune}).

\begin{itemize}[leftmargin=1.4em,itemsep=2pt]
\item A \textbf{\code{Variant}} is a named schedule: a closure \code{f(iters) -> (output, ms)} that
launches the schedule \code{iters} times after a fixed warmup and returns its output and mean
milliseconds per dispatch. The same closure both \emph{benchmarks} (many iters) and \emph{executes} (one
iter) --- this is the crate's established \code{*\_bench} shape, so no separate ``run'' and ``time'' paths
can disagree.
\item A \textbf{\code{Tuned}} is the fluent builder: collect the variants for one
\code{(op, shape-key)}, then \code{run}. It reads like the fusion builder --- fluent, no macro, no
ceremony.
\item A \textbf{\code{Tuner}} is the cache: an in-memory \code{(device, op, key) -> winner} map backed by
a per-device on-disk file. It is a \emph{plain value} constructed over an explicit cache root, so it is
fully unit-testable offline --- no GPU, no environment, no global state. The process-wide default roots
at the XDG cache directory.
\end{itemize}

\begin{lstlisting}[language=Rust,caption={The production surface. One source per op, autotuned per device,
no forks. On a cache hit only the winner launches (once, to produce its output); the losers are never
run.},label={lst:tune}]
let y = Tuned::new("rms_norm", format!("rows={rows},n={n}"))
    .variant("b64_r1",  |it| rms_norm_tuned_bench(&c, x, w, rows, n, eps, 1,  64, it))
    .variant("b128_r1", |it| rms_norm_tuned_bench(&c, x, w, rows, n, eps, 1, 128, it))
    .variant("b256_r1", |it| rms_norm_tuned_bench(&c, x, w, rows, n, eps, 1, 256, it))
    .variant("b64_r2",  |it| rms_norm_tuned_bench(&c, x, w, rows, n, eps, 2,  64, it))
    .run(&c);  // times each once per (device,op,shape), caches the winner, thereafter runs it
\end{lstlisting}

The selection logic is deliberately small. On a cache \emph{miss}, every variant is timed
\code{TUNE\_ITERS} times (25), the fastest is kept, its winning name is recorded, and its
already-computed output is returned --- no redundant re-run. On a cache \emph{hit}, only the winner is
launched, once, to produce its output; the losers are never launched. If a cached winner name is no
longer among the current variants (the set changed), that is treated as a miss and re-tuned rather than
trusting a stale name.

\paragraph{The cache is a value, not a database.} The on-disk format is one TSV file per device ---
\code{\$XDG\_CACHE\_HOME/hanzo-kernel/autotune/<device>.tsv} --- each line \code{op\textbackslash tkey\textbackslash twinner}.
No serde, no new dependency, no schema migration. Cross-process writes append and are idempotent (the
loader keeps the last line for a key), which is exactly how an advisory autotune cache should behave: a
hint that is safe to lose and safe to race on.

\section{Evaluation: the variant tables and the offline proof}
\label{sec:eval}

Two ops ship autotuned (Table~\ref{tab:variants}). Both variant sets are proven on the CPU runtime, which
needs no GPU to establish the two things under test: that the variants are numerically indistinguishable,
and that the select/cache machinery works.

\begin{table}[h]
\centering\small
\setlength{\tabcolsep}{5pt}
\begin{tabular}{@{}lll@{}}
\toprule
\textbf{Op} & \textbf{Knobs} & \textbf{Variants (all byte-identical)} \\
\midrule
\code{rms\_norm}      & cube dim $\{64,128,256\}\times$ unroll $\{1,2\}$        & \code{b64\_r1 b128\_r1 b256\_r1 b64\_r2} \\
\code{matvec\_q8\_dp4a} & cube dim $\{64,128,256\}\times$ vec.\ width $\{1,2,4,8\}$ & \code{b64\_v1 b64\_v4 b128\_v2 b256\_v8} \\
\bottomrule
\end{tabular}
\caption{The autotune variant sets (\hk{} 0.2.17). A knob value is a distinct compiled kernel; the tuner
caches the winner per \code{(device, op, shape)}. For \dpa{} matvec the vector-width knob is the ILP
factor (groups per unrolled step); the per-group int8 dot is accumulated in the same order for every
width, so all four are byte-identical.}
\label{tab:variants}
\end{table}

The committed tests assert, for both ops:
\begin{enumerate}[leftmargin=1.5em,itemsep=1pt]
\item \textbf{Byte-identity.} Every schedule's output has identical bits to the reference (RMSNorm:
exactly equal; \dpa{}: byte-identical across variants and within f32-reorder tolerance
$2\times10^{-3}$ of the per-element oracle). The knobs do not touch the numerics.
\item \textbf{Select and cache.} A first call is a miss that benchmarks all four variants
(\code{benched == 4}) and records the winner. A second call is a hit that benchmarks \emph{zero}
(\code{benched == 0}) and returns the winner's output. A fresh \code{Tuner} over the same directory
reloads the winner \emph{from disk} and skips timing. A distinct shape-key is an independent decision. A
shrunk variant set (winner absent) re-tunes rather than trusting a stale name.
\end{enumerate}
Timing on the CPU proves the select/cache mechanism; a GPU run additionally makes the \emph{winner}
meaningful, since on the CPU the schedules are near-equivalent and the chosen winner is
arbitrary-but-stable. The mechanism, not the specific winner, is what the offline tests pin down.

\section{Why first-party, not the engine's autotuner}

The lowering engine (CubeCL) ships an autotuner; it was studied and declined for three concrete reasons,
not out of not-invented-here. First, it is re-exported only from \code{cubecl\_runtime}, so adopting it
would force a \emph{second} engine crate into a facade architected so that only one line of one
\code{Cargo.toml} names the engine at all --- a deliberate ``values, not places'' boundary the whole DSL
is built to preserve. Second, its persistent cache lives under \code{target/} (via an MD5 checksum and
serde), not the XDG path this layer is required to use. Third, its generic-associated-type input
machinery is a large surface for what is, in the end, ``pick the fastest of $N$ named closures.'' The
concepts above --- variant, tuned set, cache keyed by device and shape --- are \emph{deliberately} the
same as the engine's; they are sized to this crate and pointed at XDG. A thin analog, not a new
philosophy.

\section{The takeaway}

The schedule was the last place a backend forked a kernel, and the fork was the last hiding place for the
divergence bug. Making the schedule a compile-time value folds that fork into the one source: the tuner
enumerates numerically-identical variants, times them per device, and caches the winner, so the engine is
autotuned and single-source at once. Speed is chosen; correctness is never spent. Combined with fusion as
composition (\cite{fuse}) and the oracle law (\cite{oracle}), the schedule stops being a place where two
implementations can disagree, because there is only ever one.

\begin{thebibliography}{9}
\small
\bibitem{umbrella} Hanzo AI. \emph{One Source, Every Backend: The hanzo-kernel DSL.} Paper 07 (umbrella).
\bibitem{fuse} Hanzo AI. \emph{Fusion Is Composition: Kernel Fusion as the Functor Law.} Paper 07.1.
\bibitem{tune} Hanzo AI. \emph{The Schedule Is a Value: Comptime-Specialized Autotuning.} Paper 07.2 (this paper).
\bibitem{island} Hanzo AI. \emph{Intrinsic Islands: An Oracle-Anchored Escape Hatch.} Paper 07.3.
\bibitem{oracle} Hanzo AI. \emph{One Oracle, Every Backend: The CPU Bit-Exactness Law.} Paper 07.4.
\bibitem{mmq} Hanzo AI. \emph{When the Product Won't Collapse: An int8 Tensor-Core MMQ Negative Result.} Paper 07.5.
\bibitem{method} Hanzo AI. \emph{Verify, Then Delete: Collapsing a Kernel Cartesian Product.} Paper 07.6.
\end{thebibliography}

\end{document}
